\begin{figure}
\centering
\begin{tikzpicture}[x=1cm,y=1cm]
  \node[ham/box,text width=5.3cm] (hamRoot) at (5,0)
  {Inspect the\\chain endpoints};
  \node[ham/box,text width=3.25cm] (hamSelf) at (.6,-1.65)
  {Endpoint paired\\within its own chain};
  \node[ham/box,text width=3.25cm] (hamTogether) at (5,-1.65)
  {Endpoints paired\\with each other};
  \node[ham/box,text width=3.25cm] (hamDistinct) at (9.4,-1.65)
  {Distinct cross-pairs};
  \draw[ham/arrow] (hamRoot)--(hamSelf);
  \draw[ham/arrow] (hamRoot)--(hamTogether);
  \draw[ham/arrow] (hamRoot)--(hamDistinct);
  \node[ham/box,text width=3.25cm] (hamHalf) at (.6,-3.45)
  {Regional estimate\\on the full domain};
  \node[ham/box,text width=3.25cm] (hamCS) at (5,-3.45)
  {Covariance bound\\with both origins};
  \node[ham/box,text width=3.25cm] (hamRegions) at (9.4,-3.45)
  {$s_a\le s_q$: remove $v$\\$s_a>s_q$: remove $u$};
  \draw[ham/arrow] (hamSelf)--(hamHalf);
  \draw[ham/arrow] (hamTogether)--(hamCS);
  \draw[ham/arrow] (hamDistinct)--(hamRegions);
  \node[ham/note,text width=10.5cm] at (5,-4.75)
  {The two cross-pair regions give $P/2$ and $P'/2$,\\
  with possibly different assignments: $(P+P')/2\le\max(P,P')$.};
  \draw[hamblue!35] (-1.1,-5.4)--(11.1,-5.4);
  \node[ham/note] at (1.7,-5.95) {On $s_a\le s_q$, remove $v$};
  \node[ham/dot] (hamBefore) at (-.4,-6.8) {};
  \node[ham/dot] (hamA) at (1.4,-6.8) {};
  \node[ham/dot] (hamB) at (3.2,-6.8) {};
  \node[ham/dot] (hamQ) at (1.4,-8.0) {};
  \node[ham/dot] (hamV) at (3.2,-8.0) {};
  \draw[ham/arrow] (hamBefore)--node[above,font=\small] {$G_a$} (hamA);
  \draw[ham/arrow] (hamA)--node[above,font=\small] {$G_b$} (hamB);
  \draw[ham/arrow] (hamQ)--(hamV);
  \draw[ham/pair] (hamA)--(hamV);
  \node[font=\small,above=2pt] at (hamA) {$a$};
  \node[font=\small,above=2pt] at (hamB) {$b$};
  \node[font=\small,below=2pt] at (hamQ) {$q$};
  \node[font=\small,below=2pt] at (hamV) {$v$};
  \draw[ham/arrow] (4.2,-7.5)--(5.3,-7.5);
  \node[font=\small] at (4.75,-7.0) {$N_{1,e_0}(G_v)$};
  \node[ham/dot] (hamAfterBefore) at (6.4,-6.8) {};
  \node[ham/dot] (hamAfterB) at (9.7,-6.8) {};
  \node[ham/dot] at (8.05,-8.0) {};
  \draw[ham/arrow] (hamAfterBefore)--node[above,font=\small] {$G_a*G_b$} (hamAfterB);
  \node[font=\small,above=2pt] at (hamAfterB) {$b$};
  \node[font=\small,below=2pt] at (8.05,-8.0) {$q$};
  \node[ham/note,text width=5.2cm] at (8.05,-9.5)
  {Integrate out $v$ and $a$;\\keep the other predecessors unchanged.};
\end{tikzpicture}
\caption[Endpoint elimination in two ordered chains]{The endpoint cases in the finite two-chain product bound. The empty
matching is the induction base; if only one chain is nonempty, use the
within-chain case.
Here $u,v$ are the chain endpoints, $a$ is paired with $v$,
and $q$ precedes $v$ in the distinct-cross-pair case. The lower sketch
shows one local fragment used in the regional fusion step, with
$b$ the immediate successor of $a$; other pairings are suppressed. Deleting $v$ on $s_a\le s_q$ extracts
$N_{1,e_0}(G_v)$, where $e_0=\{a,v\}$. If the unpaired $a$ is terminal,
integrate it to $N_0(G_a)$ instead; otherwise fuse its two consecutive
kernels by one-sided spacetime convolution. The mutually paired endpoint case retains
both predecessor origins, including their time difference. The argument
chooses a finite maximum of products; it does not produce one assignment
valid for all kernel families. See
\ProofRef{Lemma}{L:pairing-product-envelope}.}
\label{Fig:endpoint-elimination}
\end{figure}
